\begin{section}{Relations}

A \textbf{relation} $\sim$ on a set $S$ defines a subset $R$ of $S \times S$, and vice versa.
\[
R \coloneqq \setbuilder*{(x, y) \in S \times S}{x \sim y}, \quad x \sim y \iff (x, y) \in R.
\]

\bigskip

\begin{definition} \label{def:EquivalenceRelation}
    A relation $\sim$ on a set $S$ is called an \textbf{equivalence relation} if it satisfies the following properties:
    \begin{itemize}
        \item Reflexivity: $x \sim x$ for all $x \in S$.
        \item Symmetry: If $x \sim y$, then $y \sim x$ for all $x, y \in S$.
        \item Transitivity: If $x \sim y$ and $y \sim z$, then $x \sim z$ for all $x, y, z \in S$.
    \end{itemize}
\end{definition}

\bigskip

\begin{definition} \label{def:Partition}
    A \textbf{partition} of a set $S$ is a collection $\set{A_{\alpha}}_{\alpha \in I}$ of nonempty subsets of $S$ such that
    \begin{itemize}
        \item $A_{\alpha} \cap A_{\beta} = \emptyset$ for all $\alpha, \beta \in I$ with $\alpha \neq \beta$.
        \item $\displaystyle\bigcup_{\alpha \in I} A_{\alpha} = S$.
    \end{itemize}
\end{definition}

\bigskip

\begin{definition} \label{def:EquivalenceClass}
    Let $\sim$ be an equivalence relation on a set $S$. For each $x \in S$,
    the \textbf{equivalence class} of $x$ is defined as
    \[
    [x] \coloneqq \setbuilder*{y \in S}{y \sim x}.
    \]
    We denote the set of all equivalence classes of $S$ by $S / \sim \; \coloneqq \setbuilder*{[x]}{x \in S}$.
\end{definition}

\bigskip

\begin{theorem} \label{thm:EquivalenceClassesFormAPartition}
    Let $\sim$ be an equivalence relation on a set $S$. For any $x, y \in S$, we have either
    \[
    [x] = [y] \quad \text{or} \quad [x] \cap [y] = \emptyset.
    \]
    Therefore, $S / \sim$ is a partition of $S$.
\end{theorem}

\begin{proof}
    Suppose $[x] \cap [y] \neq \emptyset$. Then there exists $z \in S$ such that $z \in [x]$ and $z \in [y]$.
    This implies that $z \sim x$ and $z \sim y$. By the symmetry and transitivity of $\sim$,
    we have $x \sim z$ and $z \sim y$, which implies that $x \sim y$.
    \\[3mm]
    Therefore, for any $w \in [x]$, we have $w \sim x$ and $x \sim y$, which implies that
    $w \sim y$, so $w \in [y]$. This shows that $[x] \subseteq [y]$. Similarly, we can show that $[y] \subseteq [x]$,
    so $[x] = [y]$.
\end{proof}

\bigskip

\begin{theorem} \label{thm:PartitionDefinesEquivalenceRelation}
    Let $\set{A_{\alpha}}_{\alpha \in I}$ be a partition of a set $S$. Define a relation $\sim$ on $S$ by
    \[
    x \sim y \iff x, y \in A_{\alpha} \text{ for some } \alpha \in I.
    \]
    Then $\sim$ is an equivalence relation on $S$,
    and the equivalence classes of $\sim$ are precisely the sets $A_{\alpha}$.
\end{theorem}

\begin{proof}
    We need to show that $\sim$ is reflexive, symmetric, and transitive.
    \begin{itemize}
        \item Reflexivity: For any $x \in S$, there exists $\alpha \in I$ such that $x \in A_{\alpha}$, so $x \sim x$.
        \item Symmetry: If $x \sim y$, then there exists $\alpha \in I$ such that $x, y \in A_{\alpha}$, so $y \sim x$.
        \item Transitivity: If $x \sim y$ and $y \sim z$, then there exist $\alpha, \beta \in I$ such that
        $x, y \in A_{\alpha}$ and $y, z \in A_{\beta}$. Since $y \in A_{\alpha} \cap A_{\beta}$ and the sets
        $\set{A_{\alpha}}_{\alpha \in I}$ are disjoint, we must have $\alpha = \beta$, so $x, z \in A_{\alpha}$ and $x \sim z$.
    \end{itemize}
    Therefore, $\sim$ is an equivalence relation on $S$. The equivalence class of any $x \in S$
    is precisely the set $A_{\alpha}$ containing $x$, so the equivalence classes of $\sim$ are exactly the sets $A_{\alpha}$.
\end{proof}

\end{section}